\section{Представление деревьев в ЭВМ}
\label{sec:tree-representation}

\subsection{Общие подходы}

Дерево можно хранить теми же способами, что и произвольный граф: матрицей или списками смежности. Однако наличие корня, единственного родителя и иерархической структуры позволяет использовать более компактные представления.

Выбор структуры зависит от требуемых операций:
\begin{itemize}
    \item переход к родителю;
    \item перебор детей;
    \item вставка и удаление поддерева;
    \item поиск по ключу;
    \item последовательное хранение и сериализация.
\end{itemize}

\subsection{Массив родителей}

Для корневого дерева с вершинами $1,\ldots,n$ хранится массив
\[
    p[1\ldots n],
\]
где $p[v]$ --- номер родителя вершины $v$, а для корня используется специальное значение, например $0$ или $-1$.

Память составляет $O(n)$. Переход к родителю выполняется за $O(1)$, но поиск всех детей вершины без дополнительного индекса требует $O(n)$. Поэтому массив родителей удобен, когда движение происходит преимущественно снизу вверх.

Глубину вершины можно находить последовательным переходом по родителям. Для многократных запросов применяют двоичный подъём: дополнительно хранят $up[v][k]$ --- предка вершины $v$ на расстоянии $2^k$. Тогда предок на заданной высоте и наименьший общий предок вычисляются за $O(\log n)$ (представляем заданную высоту как степени двойки, тогда <<прыжков>> будет столько, сколько степеней двойки -- порядка $\log n$).

\subsection{Списки детей}

Для каждой вершины хранится список её детей. Общий объём памяти равен $O(n)$, поскольку в дереве $n-1$ ребро. Перебор детей вершины занимает время, пропорциональное их числу.

Рассмотрим для примера дерево:
\[
0 \to \{1,2,3\}, \qquad 3 \to \{4,5\},
\]
где вершины $1,2,4,5$ являются листьями.

На практике используют:
\begin{itemize}
	\item динамические списки или векторы детей. Для приведённого дерева можно хранить
	\[
	children[0]=\{1,2,3\},\qquad
	children[3]=\{4,5\},
	\]
	а списки остальных вершин оставить пустыми;
	
	\item один массив детей и массив начал диапазонов, если дерево неизменно. Например,
	\[
	children=\{1,2,3,4,5\},
	\]
	\[
	start=\{0,3,3,3,5,5,5\}.
	\]
	Тогда дети вершины $v$ находятся в диапазоне
	\[
	children[start[v]\,..\,start[v+1]).
	\]
	Например, для вершины $0$ это элементы с индексами $0,1,2$, то есть
	$\{1,2,3\}$, а для вершины $3$ --- элементы с индексами $3,4$, то есть
	$\{4,5\}$;
	
	\item связное представление ``первый ребёнок --- следующий брат''.
\end{itemize}

В последнем случае каждый узел содержит две ссылки:
\[
firstChild(v),\qquad nextSibling(v).
\]
Для того же примера можно задать
\[
firstChild(0)=1,\qquad
nextSibling(1)=2,\qquad
nextSibling(2)=3,
\]
а для вершины $3$:
\[
firstChild(3)=4,\qquad
nextSibling(4)=5.
\]
Отсутствующие дети и следующие братья обозначаются нулевым указателем
(или специальным значением).

Так любое дерево произвольной степени представляется структурой, похожей на бинарное дерево, с постоянным числом ссылок на узел.

\subsection{Связное представление бинарного дерева}

Узел бинарного дерева обычно содержит данные и указатели на левое и правое поддеревья:
\begin{verbatim}
struct Node {
    Data value;
    Node* left;
    Node* right;
    Node* parent; // необязательно
};
\end{verbatim}

Память --- $O(n)$. Переход к ребёнку выполняется за $O(1)$. Наличие указателя на родителя упрощает подъём вверх, но увеличивает память. В языках с автоматическим управлением памятью вместо указателей используются ссылки.

\subsection{Последовательное представление бинарного дерева}

Полное или почти полное бинарное дерево удобно хранить в массиве. При нумерации с единицы для вершины $i$:
\[
    left(i)=2i,
    \qquad
    right(i)=2i+1,
    \qquad
    parent(i)=\left\lfloor\frac{i}{2}\right\rfloor.
\]
При нумерации с нуля используются индексы $2i+1$, $2i+2$ и $\lfloor(i-1)/2\rfloor$ соответственно.

Так хранят бинарные кучи. Представление не требует указателей. Для сильно разреженного дерева оно неэффективно, поскольку возникает много пустых позиций.

\subsection{Обходы как способ кодирования (необязательно)}

Для бинарного дерева определены три основных обхода:
\begin{itemize}
    \item прямой: корень, левое поддерево, правое поддерево;
    \item симметричный: левое поддерево, корень, правое поддерево;
    \item обратный: левое поддерево, правое поддерево, корень.
\end{itemize}

Один обход в общем случае не определяет дерево однозначно. Для восстановления бинарного дерева с различными ключами достаточно пары обходов ``прямой + симметричный'' или ``обратный + симметричный''. При сериализации можно хранить прямой обход вместе с маркерами пустых поддеревьев; тогда представление однозначно.

Для упорядоченного дерева используется скобочная запись. Например, вершина и её дети записываются рекурсивно как
\[
    v(T_1,T_2,\ldots,T_k).
\]

\subsection{Код Прюфера}

Свободное помеченное дерево на $n$ вершинах кодируется последовательностью длины $n-2$. Алгоритм построения:
\begin{enumerate}
    \item выбрать лист (висячую вершину) с наименьшей меткой;
    \item записать метку его единственного соседа;
    \item удалить лист и повторить действия до двух оставшихся вершин.
\end{enumerate}

Код Прюфера занимает $O(n)$ памяти и однозначно определяет дерево. Число вхождений метки $v$ в код равно $d(v)-1$. Построение и декодирование можно выполнить за $O(n\log n)$ с очередью с приоритетом или за $O(n)$ при специальной организации данных.

\subsection{Польская запись деревьев выражений}

Дерево арифметического выражения можно кодировать префиксной или постфиксной записью. Например, выражению
\[
    (a+b)c
\]
соответствуют:
\[
    \text{префиксная запись } *\,+\,a\,b\,c,
    \qquad
    \text{постфиксная запись } a\,b\,+\,c\,*.
\]
Постфиксная запись вычисляется стеком без скобок и удобна для компиляторов.

\subsection{Что важно сказать на экзамене}

Для корневого дерева основные структуры --- массив родителей и списки детей; для бинарного --- узлы с двумя ссылками или хранение в массиве. Нужно сравнить память и стоимость переходов. По желанию упомянуть сериализацию обходами и код Прюфера.

\newpage